\chapter*{Abstract} %*-Variante sorgt dafür, das Abstract nicht im Inhaltsverzeichnis auftaucht
\phantomsection\pdfbookmark{Abstract}{abstract}
Under the Regulation (EU) 2023/1230, cybersecurity has become a mandatory part of machinery conformity assessment. Cybersecurity-related properties of industrial gateways serve as evidence within this assessment, especially when these gateways influence safety-related functions. Since established functional validation methods currently lack verification of security functions, this thesis develops and validates a testing methodology to address this gap.

Following the \gls{acr:DSRM} methodology, fourteen component-level security requirements were extracted from Annex III obligations and mapped via the draft standard prEN 50742 to the \gls{acr:IEC} 62443-4-2 component requirements. Each requirement was linked to affected assets, a \gls{acr:STRIDE} threat, and a corresponding security function within a two-level test catalogue. An existing test framework of ifm was integrated through a four-outcome applicability assessment and classified using the \gls{acr:IEC} 62443-4-1 (\gls{acr:SVV}) verification taxonomy. Demonstration occurred on a Safe Remote \gls{acr:I/O} device under the direct-access network position \gls{acr:NV1}.

During evaluation, fourteen requirements, nine assets, and ten company tests were assessed: five were applicable and executed, three were not applicable, one was blocked, and one original procedure was not applicable although the tool itself produced a usable transport-level observation. A set of evidence-backed test activities was performed under \gls{acr:NV1}, mostly supported by stored artefacts. Two requirements were assessed as Fail due to evidenced capability deviations, while the remaining twelve were Not Assessed because their complete acceptance conditions require the passive-observation or in-path network variants, an induced software or configuration change, or the engineering tool.

The central methodological finding demonstrates that a safe-state response and cybersecurity conformity are separable properties requiring independent verification. A defined safe state after a fault does not inherently demonstrate accountability, integrity, or traceability.



\clearpage

\chapter*{Kurzfassung} %*-Variante sorgt dafür, das Abstract nicht im Inhaltsverzeichnis auftaucht
\phantomsection\pdfbookmark{Kurzfassung}{kurzfassung}

Gemäß der Verordnung (EU) 2023/1230 ist die Cybersicherheit zu einem obligatorischen Bestandteil der Konformitätsbewertung von Maschinen geworden. Cybersicherheitsrelevante Eigenschaften industrieller Gateways dienen als Nachweis im Rahmen dieser Bewertung, insbesondere wenn diese Gateways sicherheitsrelevante Funktionen beeinflussen. Da es bei der etablierten funktionalen Validierung derzeit an Überprüfung von Sicherheitsfunktionen mangelt, entwickelt und validiert diese Arbeit eine Testmethodik, um diese Lücke zu schließen.


In Anlehnung an die \gls{acr:DSRM}-Methodik wurden vierzehn Sicherheitsanforderungen auf Komponentenebene aus den Verpflichtungen in Anhang III extrahiert und über den Normenentwurf prEN 50742 den Komponentenanforderungen der \gls{acr:IEC} 62443-4-2 zugeordnet. Jede Anforderung wurde mit den betroffenen Assets, einer \gls{acr:STRIDE}-Bedrohung und einer entsprechenden Sicherheitsfunktion innerhalb eines zweistufigen Testkatalogs verknüpft. Ein bestehendes Test-Framework der ifm wurde durch eine Anwendbarkeitsbewertung mit vier Ergebnissen integriert und anhand der Verifikationstaxonomie der Norm \gls{acr:IEC} 62443-4-1 (\gls{acr:SVV}) klassifiziert. Die Demonstration erfolgte auf einem sicheren Remote-\gls{acr:I/O}-Gerät unter der Netzwerkposition mit direktem Zugriff \gls{acr:NV1}.

Im Rahmen der Bewertung wurden vierzehn Anforderungen, neun Assets und zehn Unternehmenstests geprüft: Fünf waren anwendbar und wurden ausgeführt, drei waren nicht anwendbar, eine wurde blockiert, und bei einem ursprünglichen Verfahren war die Anwendbarkeit zu verneinen, während das Werkzeug selbst eine verwertbare Beobachtung auf Transportebene lieferte. Eine Reihe evidenzbasierter Testaktivitäten wurde unter \gls{acr:NV1} durchgeführt und größtenteils durch gespeicherte Artefakte belegt. Zwei Anforderungen wurden aufgrund nachgewiesener Fähigkeitsabweichungen als nicht erfüllt (Fail) bewertet, während die verbleibenden zwölf als nicht bewertet (Not Assessed) eingestuft wurden, da ihre vollständige Erfüllung die passive Beobachtungs- oder Inline-Netzwerkvariante, eine induzierte Software- oder Konfigurationsänderung oder das Engineering-Tool erfordert hätte.

Die zentrale methodische Erkenntnis zeigt, dass eine Reaktion im sicheren Zustand und die Konformität mit den Cybersicherheitsanforderungen trennbare Eigenschaften sind, die einer unabhängigen Verifizierung bedürfen. Ein definierter sicherer Zustand nach einem Fehler belegt nicht automatisch Rechenschaftspflicht, Integrität oder Rückverfolgbarkeit.





\clearpage